%!TEX root=../../main.tex

In 2012, AlexNet won the ImageNet Large Scale Visual Recognition Challenge
\citep{deng2009imagenet} and completed the ascendency of deep neural networks in
machine learning, a process which (arguably) began with Geoffrey Hinton's
famous MNIST digit generation demo in 2006 \citep{hinton2006deep}.
AlexNet won the competition with a top-five accuracy of \( 84.7\%  \) and at
the time was considered a ``large, deep convolutional neural network'' based on
its model size of sixty million parameters \citep{russakovksy2015imagenet}.
In comparison, the recent foundation model Florence
\citep{yuan2021florence} achieves top-five accuracy greater than \( 99\% \)
using 893 million parameters.
In addition to being much larger than AlexNet, Florence uses modern techniques
like contrastive learning, image transformers, and extensive pre-training on
800 million additional images.
The gaps between AlexNet and Florence in modeling, training data, and test
performance reflect the rapid advances machine learning has made since 2012.

While AlexNet took approximately seven days to train on two NVIDIA GTX 580 GPUs
\citep{krizhevksy2012alexnet}, Florence required 10 days to pre-train on 512
NVIDIA A100 GPUs.
This means Florence used roughly 365 times more GPU hours on machines capable
of either 12 (32 bit precision) or 48 (16 bit precision)
times more FLOPS and with 20 times more memory than those used for AlexNet.
Accounting for Florence's size and additional data yields a training
efficiency between \( 0.68 \) and \( 2.74 \) relative to AlexNet
\citep{a100, gpucomparison}.%
\footnote{Florence uses mixed-precision training, so an exact performance
    comparison is difficult.}
The point of these back-of-the-envelope calculations is to expose a dominant
trend in machine learning:
increasingly large models and data sets have led to remarkable improvements in
test performance, but training costs have increased proportionally and
training efficiency has stagnated.
In other words, advances in optimization for neural networks have not matched
the advances in neural networks themselves.

The key challenge to designing better optimizers for neural networks is
non-convexity.
Unlike matrix factorization and other benign non-convex problems
\citep{burer2005low}, neural networks admit many saddle points and
local extrema \citep{safran2018local, dauphin2014saddles}.
As a result, global optimization for neural networks is NP-hard in general
\citep{froese2023hard, blum1988complete}.
Faced with this formal hardness guarantee, practitioners rely on stochastic
gradient methods (SGMs) to approximately minimize the training objective.
However, SGMs do not guarantee convergence to a local optimum \citep{ge2015escaping}
and existing methods rarely certify convergence to a stationary point of any
type \citep{goodfellow2016deeplearning}.
They are also slow from a theoretical perspective:
lower bounds show that SGMs require at least \( \Omega(\epsilon^{-4}) \) iterations
to compute an \( \epsilon \)-stationary point \citep{arjevani2023lower, drori2020complexity}.
In stark contrast, for smooth, convex functions a stationary point can be
computed in \( \tilde{O}(\epsilon^{-1/2}) \) iterations using acceleration
\citep{nesterov2018lectures}.

In addition to being theoretically slow, SGMs are also highly sensitive to
their hyper-parameters.
For example, they converge slowly or may even diverge depending on the choice
of step-size and/or momentum parameters \citep{zhang2022adam}.
Parameters like the random seed complicate replications and can produce model
churn, where networks learned using the same procedure give different
predictions for the same inputs \citep{henderson2018deep,
    bhojanapalli2021reproducibility}.
Although some authors have argued that these hyper-parameters make
SGMs more adaptable \citep{choi2019empirical}, attaining this adaptivity
requires the parameters to be carefully tuned for each problem
\citep{sivaprasad2020tuning}.
Non-convexity is again a critical challenge:
standard solutions to tuning like line-search
\citep{armijo1966ls} and restarts \citep{odonoghue2015restarts} cannot be
applied to non-smooth, non-convex neural networks.
Instead, practitioners are forced to rely on domain knowledge and costly,
input-intensive algorithms like grid-search.

\begin{figure}[t]
    \centering
    \input{figures/chapter_one/non_convex.tex}
    \caption{Illustration of critical points for a non-convex function.
        Although \( \tilde{x}_1, \tilde{x}_2 \), and \( x^* \) are all
        stationary points of \( f \), only \( x^* \) is a global minimum.
        In comparison, \( \tilde{x}_1 \) is a local minimum and 
        \( \tilde{x}_2 \) is a local maximum.
        Characterizing global minima is challenging without stronger
        optimality conditions.
    }
    \label{fig:non-convex}
\end{figure}

Non-convexity is also the main barrier to understanding basic properties of the
neural network training problem.
For non-convex problems, stationarity of the training objective does not imply
global optimality of the network weights (see \cref{fig:non-convex}).
Without strong optimality conditions, characterizations of global optima have
been limited to embedding principles
\citep{fukumizu2000plateau, zhang2021embedding}, mode connectivity
\citep{garipov2018mode, draxler2018barriers}, and existence
arguments for optimal manifolds under restricted architectures
\citep{cooper2021global}.
This is in stark contrast to generalized linear models with \( \ell_2 \)
or \( \ell_1 \) penalties;
for example, it is well-known that the lasso \citep{tibshirani1996regression}
has a piece-wise linear regularization path \citep{osborne2000new, efron2004least},
a polyhedral solution set \citep{tibshirani2013unique}, and
admits efficient algorithms for computing
minimal solutions \citep{tibshirani2013unique}.
Optimization has not just fallen behind on training algorithms for neural
networks, but also on the fundamental theory of neural network training
problems.

This thesis tackles the optimization gap for neural networks using convex
reformulations.
Convex reformulations re-write the non-convex training problem for a neural
network as a convex program by lifting it into a higher dimensional space
\citep{pilanci2020convexnn}.
Crucially, the resulting problem is finite-dimensional.
As such, convex reformulations are distinct from convexifications of neural
networks based on infinite-width relaxations
\citep{bach2017breaking, bengio2006convex}.
Indeed, convex reformulations can be viewed as providing a finite
discretization for the measure optimization problem arising in the
infinite-width setting.
Since their introduction for two-layer ReLU networks by
\citet{pilanci2020convexnn}, convex reformulations have been developed for many
other shallow architectures, including for convolutions
\citep{sahiner2020convex,ergen2021implicit,gupta2021exact}, vector-outputs
\citep{sahiner2021vector}, batch normalization \citep{ergen2021demystifying},
generative models \citep{sahiner2021hidden} and polynomial activations
\citep{bartan2025polynomial}.

We focus on convex reformulations of two-layer neural networks with ReLU
activation functions.
From a theoretical perspective, two-layer ReLU networks are highly attractive
because they are piece-wise linear predictors as well as universal
approximators \citep{huang2020universal}.
However, ReLU activations are also among the most popular in machine learning
\citep{szandala2020review} and two-layer ReLU networks are fundamental building
blocks of many important models.
For instance, deep ReLU networks can be modeled as stacks of two-layer modules
and training by greedy sequential optimization of these blocks is competitive
with end-to-end training \citep{belilovsky2019greedy}.
Two-layer ReLU networks are basic components of the standard transformer block
and can be used to model self attention \citep{sahiner2022attention}.
On their own, two-layer ReLU networks are widely used for reinforcement
learning \citep{gaur2023global}, prediction on edge devices
\citep{teerapittayanon2017distributed}, and word-embeddings
\citep{mikolov2013word}.
As such, advances in training and theory for two-layer networks stand to
impact many areas of machine learning.

We start by leveraging convex reformulations to develop fast, global
optimization methods for two-layer ReLU networks.
Access to an equivalent convex program allows us to side-step optimization
lower-bounds for non-convex, non-smooth functions and develop accelerated
solvers based on the FISTA algorithm \citep{beck2009fista} and on
augmented Lagrangian methods \citep{rockafellar1976augmented}.
Key to this approach is the development of an unconstrained relaxation of the
convex ReLU problem called a gated ReLU network \citep{fiat2019decoupling}.
We show that gated ReLU networks with non-singular activation patterns form the
same hypothesis space as ReLU networks and provide approximation bounds for
regularized training problems.
Experimentally, our convex solvers are orders of magnitude faster than SGMs
while requiring with little-to-no hyper-parameter tuning.

Next, we use convex reformulations to characterize the set of global
solutions to the non-convex training problem for two-layer ReLU networks.
Our work reveals the existence of a non-trivial manifold of global optima;
in the convex parameterization, this manifold corresponds to a convex polytope
living in an affine subspace.
We show that vertices of the polytope correspond to highly neuron-sparse
networks and derive a neuron pruning algorithm to compute these minimal models
To the best of our knowledge, this is the first optimal neuron pruning
algorithm for neural networks.
We also prove many other continuity and sensitivity results as well as extend
our characterization of the optimal set to include all critical points of the
training objective.

Since these optimality results address only two-layer networks, we conclude by
developing convex reformulations for arbitrarily deep ReLU networks.
Our approach is based on a novel layer elimination lemma, which we use to
to show that ReLU networks are equivalent to tensor optimization problems with
non-convex, low-rank decomposition constraints.
For sufficiently wide networks (although still finite), these low-rank
constraints vanish and we obtain a fully convex tensor program with a
recursive cone constraint.
Similar to the two-layer setting, we develop a unconstrained Gated ReLU
relaxation and prove equivalence of the model spaces for non-singular
activation patterns.
Finally, we study both deep ReLU and deep gated ReLU models under a variety
of regularization schemes.

The rest of the thesis is organized as follows:
In \cref{chapter:fast-optimization}, we develop the theory of gated ReLU
networks and present our convex solvers for two-layer models.
Next, \cref{chapter:solution-functions} derives the structure of the optimal
set for two-layer ReLU networks and establishes many results around pruning,
sampling, continuity, and sensitivity.
Finally, \cref{chapter:deep-reformulations} extends convex reformulations
to deep, fully-connected ReLU networks.
All proofs are postponed to the appropriate section of the appendix.

